%============================================================================%
%
%
%	DOCUMENMT DEFINITION
%
%
%============================================================================%

% we use report class for thesis design
% 11pt font is way bettern to read than 12pt
% we want to make a title page
% two side allows us to make the page layout align by even and odd
% openright tells the compiler that the first page is a right page
\documentclass[pdftex,11pt,titlepage,twoside,openright]{report}	


%----------------------------------------------------------------------------------------
%	ENCODING
%----------------------------------------------------------------------------------------

% for suporting multi platform
\usepackage[utf8x]{inputenc} 	

% natbib is what you want for bibliography
\usepackage[square,authoryear]{natbib}


%----------------------------------------------------------------------------------------
%	ALL DECLARATIONS
%----------------------------------------------------------------------------------------

% here are all our declarations 
%----------------------------------------------------------------------------------------
%	COLOR DEFINITIONS
%----------------------------------------------------------------------------------------

% we use the xtended color package
\usepackage[table,usenames,dvipsnames]{xcolor}

% citation reference color
\definecolor{CITATION_COL}{RGB}{121,0,226}

% urls color
\definecolor{SECTION_COL}{RGB}{0, 117,226}

% internal references color
\definecolor{LINK_COL}{rgb}{0,0,0}

% main color
\definecolor{MAIN_COL}{RGB}{191, 87, 0}
% code highlight color
\definecolor{CODE_COL}{rgb}{0.8,0.8,0.8}

% some black fallback
\definecolor{DEF_COL}{RGB}{226,109,0}

%----------------------------------------------------------------------------------------
%	REFERENCES URL CITATION STYLES AND SO ON
%----------------------------------------------------------------------------------------

% apply color definitions to hyperref package
\usepackage[colorlinks=true,citecolor=CITATION_COL,urlcolor=SECTION_COL,linkcolor=LINK_COL]{hyperref}

% make urls directly open the browser onclick
\usepackage{url}
\urlstyle{same}


%----------------------------------------------------------------------------------------
%	FONT DEFINITIONS
%----------------------------------------------------------------------------------------

% better typography rendering scheme
\usepackage[protrusion=true,expansion=true]{microtype} 

 % we use the standard Times font
\usepackage{pslatex}

 % required for accented characters
\usepackage[T1]{fontenc}

% change line spacing here, readability 
% benefits from a slight increase by default
\linespread{1.25} 

% allows more font size definitions
\usepackage{moresize}

%----------------------------------------------------------------------------------------
%	GEOMETRY  DEFINITIONS
%----------------------------------------------------------------------------------------

% define page styles using geometry
\usepackage[bottom=3cm,marginparwidth=1.5cm, marginparsep=0.5cm]{geometry}

 %define A4 Paper
\geometry{a4paper}	

% define margins for even and odd pages
\setlength{\oddsidemargin}{15.5pt}
\setlength{\evensidemargin}{15.5pt}

% paragraph indent and skip
\setlength{\parindent}{0mm}
\setlength{\parskip}{1mm }

% short command for colorizing a text as a note
% not used within this thesis anymore
\newcommand{\NOTE}[1]{\textit{\textcolor{DEF_COL}{#1}}}
	%\marginpar{\sffamily \MakeUppercase{\tiny{{#1}}}}


%----------------------------------------------------------------------------------------
%	GRAPHICS  DEFINITIONS
%----------------------------------------------------------------------------------------

% required for including pictures
\usepackage[pdftex]{graphicx}

 % make it possible to include more than one captioned figure/table in a single float 
\usepackage{subfig}

% make it possible to control text
% positioning around figures
\usepackage{wrapfig}

 % for much better looking tables
\usepackage{booktabs}	

% for better arrays (eg matrices) in maths		
\usepackage{array} 


%----------------------------------------------------------------------------------------
%	ENVIRONMENT  DEFINITIONS
%----------------------------------------------------------------------------------------

 % very flexible & customisable lists
%(eg. enumerate/itemize, etc.)				
\usepackage{paralist}				

% adds environment for commenting
% out blocks of text & for better verbatim
\usepackage{verbatim}				

% allows multiple rows to be merged to one row
\usepackage{multirow}

%----------------------------------------------------------------------------------------
% CUSTOM STRUT FOR EMPTY BOXES
%----------------------------------------- -----------------------------------------------

% when boxes become a pain this might help
\newcommand{\mystrut}{\rule[-.3\baselineskip]{0pt}{\baselineskip}}

%----------------------------------------------------------------------------------------
%	HEADER / FOOTER  DEFINITIONS
%----------------------------------------------------------------------------------------

% header definitions package
 % this should be set AFTER setting up the page geometry
\usepackage{fancyhdr}

% we want to shorten the display of the chapter in the header
\renewcommand{\chaptermark}[1]
{
  \markboth{#1}{}
}

% but we want to show extended info of the current section
\renewcommand{\sectionmark}[1]{\markright{\thesection\ #1}}

% we can even manipulate the font family
%\lhead{\sffamily Chapter \thechapter}

% this defines, what content should be displayed in header and footer
% and on which side (left-right) as well as on even and odd pages
\fancyhf[FLE,FRO]{}
\fancyhf[HLE,HRO]{\colorbox{MAIN_COL}{\sffamily\LARGE\textcolor{white}{\thepage}}}
\fancyhf[HRE,HLO]{\sffamily \textcolor{MAIN_COL}{\leftmark}}
\cfoot{}

%override plain page style for plain pages
\fancypagestyle{plain}{ %
  \fancyhf{} % remove everything
  \renewcommand{\headrulewidth}{0pt} % remove lines as well
  \renewcommand{\footrulewidth}{0pt}
  \fancyhf[HLE,HRO]{\colorbox{MAIN_COL}{\sffamily\large\textcolor{white}{\thepage}}}
  \cfoot{}
}

%----------------------------------------------------------------------------------------
%	SECTION TITLE APPEARANCE
%----------------------------------------------------------------------------------------

% we want to customize our section title 
% to make it look more appealing
\usepackage{sectsty}

% for dimensioning and other math stuff
\usepackage{amsmath}

% modify bg colors
\usepackage[some]{background}

% draw custom graphics progammatically
\usepackage{tikz}				
\usetikzlibrary{shapes, backgrounds,mindmap, trees}

% this is our custom section page layout
\newcommand{\SECTBOX}[1]
{
	\begin{tikzpicture}[remember picture,overlay]
 		\path [fill=MAIN_COL] (current page.west)rectangle (current page.north east); 
	\end{tikzpicture}
}

% this is our custom title page layout
\newcommand{\TITLEBOX}
{
	\begin{tikzpicture}[remember picture,overlay]
 		\path [fill=MAIN_COL] (current page.north west) rectangle (25.0,-3.5); 
	\end{tikzpicture}
}


 % (See the fntguide.pdf for font help)
\allsectionsfont{\sffamily\mdseries\upshape}

% for customizing the title page
\usepackage[explicit]{titlesec}

% font and color def for section
\titleformat{\section}{\LARGE\sffamily}{\textcolor{LINK_COL}{\thesection \hspace{12pt} #1}}{0pt}{}

% font and color def for subsection
\titleformat{\subsection}{\large\bf\sffamily}{\textcolor{SECTION_COL}{\thesubsection \hspace{12pt} #1}}{0pt}{}

% note: remove \thesection or \thesubsection to leave out the 
% numbering at the start of the section / subsection headline


%----------------------------------------------------------------------------------------
%	TABLE OF CONTENT APPEARANCE
%----------------------------------------------------------------------------------------

% deprecated
%\usepackage[nottoc,notlof,notlot]{tocbibind} % Put the bibliography in the ToC

% Alter the style of the Table of Contents
%\usepackage[titles,subfigure]{tocloft}

%\renewcommand{\cftsecfont}{\rmfamily\mdseries\upshape}

%\renewcommand{\cftsecpagefont}{\rmfamily\mdseries\upshape} % No bold!

% renames abstract to executive summary
\renewcommand{\abstractname}{Executive Summary} 		


%----------------------------------------------------------------------------------------
%	CUSTOM HRULE
%----------------------------------------------------------------------------------------

% a custom horizontal rule / line
\newcommand{\HRule}{\rule{\linewidth}{0.5mm}}

\newcommand{\ThinHRule}{\textcolor{CODE_COL}{\rule{0.5\linewidth}{0.05mm}}}


%----------------------------------------------------------------------------------------
%	ENVIRONMENTS
%----------------------------------------------------------------------------------------

% extended listin options
% can be used for code display
\usepackage{listings}

% syntax highlight def for listings
\lstset{ %
  	backgroundcolor=\color{CODE_COL},
	language=C,
	breaklines=true
}

% allows us to define pseudocode algorithms
\usepackage[ruled,linesnumbered]{algorithm2e}

% extended table definitions
\usepackage{tabularx}

% include external pdf pages
\usepackage{pdfpages}
\usepackage{float}


%----------------------------------------------------------------------------------------
%	Code Syntax Highlight Formatting
%----------------------------------------------------------------------------------------

\lstset{
  basicstyle=\ttfamily,
  columns=fullflexible,
  showstringspaces=false,
  commentstyle=\color{gray}\upshape
}

\lstdefinelanguage{XML}
{
  morestring=[b]",
  morestring=[s]{>}{<},
  morecomment= [s]{<!}{->},
  stringstyle=\color{black},
  identifierstyle=\color{blue},
  keywordstyle=\color{orange},
  morekeywords={xmlns, pattern,type, cells, count, required} % list your attributes here which may be highlighted
}






%----------------------------------------------------------------------------------------
%	MAKE INDEX AND GLOSSARY
%----------------------------------------------------------------------------------------


\usepackage[style=long,nonumberlist,toc,xindy,acronym,nomain]{glossaries} % nomain, if you define glossaries in a file, and you use \include{INP-00-glossary}
%\loadglsentries[main]{glossary}
% or using \input:
\include{glossary}
\makeglossaries
\usepackage{makeidx}
\makeindex



%============================================================================%
%
%	BEGIN DOCUMENT
%
%============================================================================%


\begin{document}

% before the chapters start, we use roman numbering on the pages
\setcounter{page}{1}
\pagenumbering{roman}

% print title
\begin{titlepage}


%TIKZ BACKGROUND
\TITLEBOX

%CONTENT
\begin{center}



\sffamily\textsc{\huge{\textcolor{white}{The University of Texas at Austin}}}\\[4cm]


% Upper part of the page. The '~' is needed because \\
% only works if a paragraph has started.
\includegraphics[width=0.35\textwidth]{media/faculty.png}\makebox[1.5cm]{}\includegraphics[width=0.35\textwidth]{media/unilogo.png}~\\[1.5cm]



% Title
\HRule \\[0.4cm]
{
 \huge \bfseries \sffamily Vision Transformer-Assisted Analysis of
 
Neural Image Compression and Generation  \\[0.4cm] 
}

\HRule \\[0.4cm] 

\normalfont \LARGE \sffamily Master's Thesis Report \\[1.5cm]
% \normalfont \large \sffamily DRAFT v2.0 \\[0.5cm]
\normalfont \large \sffamily Official Code Repositories: \\
\normalfont \large \sffamily \url{https://github.com/kliment-slice/thesis-code} \\[0.05cm]
\normalfont \large \sffamily \url{https://github.com/kliment-slice/thesis-latex} \\[1.0cm]

% Author and supervisor
\noindent
% \begin{minipage}{0.4\textwidth}
% \begin{flushleft} \large
% \emph{Supervisors:} \\
% Prof. Dr. Alan C. Bovik\\
% Prof. Dr. Rachel A. Ward\\
% \end{flushleft}
% \end{minipage}%
\begin{minipage}{0.4\textwidth}
\begin{center} \large
\emph{Author:}\\
Kliment Minchev\\
\end{center}
\end{minipage}

\vfill


% Bottom of the page
{\large \today}

\vfill



% \vfill
\end{center}
\end{titlepage}



%pagestyle after title
\pagestyle{fancy}

% let's print the table of content, 
% the list of figures and the list of tables
\setcounter{tocdepth}{1}
% \tableofcontents
% \listoffigures
% \listoftables
{\tableofcontents \let\cleardoublepage\clearpage \listoffigures 
\let\cleardoublepage\clearpage \listoftables \let\cleardoublepage\clearpage}

%\cleardoublepage

% EXECUTIVE SUMMARY
% here we describe roughly what this work is about

\begin{abstract}

    Vision Transformers (ViT) \citep{Bai2022AAAI} and TransGANs[2] are state-of-the-art implementations of the 
    transformer attention mechanism used in generative image tasks. The transformer architecture 
    has outperformed convolutional neural networks (CNNs) in image recognition and object detection 
    due to its ability to generalize well. The attention mechanism is capable of focusing on objects 
    anywhere on the image within a single network layer (versus the variable size convolution kernels 
    across layers). The standard two CNNs combined in a generative adversarial network (GAN), can be 
    replaced with vision transformers to create a TransGAN[2]. Such a model is shown to generalize well 
    and may be capable of handling image and video compression tasks. This thesis project investigates 
    the capacity to which a transformer and potentially a TransGAN is capable of achieving high fidelity 
    video compression in order to mimic the performance of common video codecs. 
\end{abstract}

% Back to arabic numbering

\pagestyle{fancy}

\setcounter{page}{1}
\pagenumbering{arabic}


%============================================================================%
%
%	CHAPTER 1
%
%============================================================================%

\chapter{General Tools and Techniques for Your Master Thesis}

% before the chapter starts, we can write some overview 
% this is good to summarize quickly what this chapter covers
At this point it may be good for the reader to get an overview about this chapter. What does it cover? What is it's purpose in the bigger picture of your thesis? It is like a micro-abstract and it helps your audience to keep on track. This is especially useful, when your thesis is an interdisciplinary work.



\newpage
%--
%	CHAPTER 1
%--

\section{Motivation}

Transformers are presently considered to hold a great promise for the future of Deep Learning
as a step towards Artificial General Intelligence.
Due to their architecture, they are more generalizable, less prone to overfitting, and able 
to learn highly complex representations. The Transformer architecture has already been proven 
to make obsolete Recurrent Neural Networks (RNNs) in natural language models. Furthermore, the 
Vision Transformer (ViT) has outperformed certain Convolutional Neural Networks (CNNs) in image classification tasks. \citep{dosovitskiy2020vit}

Figure 1.1 below shows an increase in the popularity of research related to Vision Transformers.

\begin{figure}[H]
	\begin{center}
	\includegraphics[width=1\textwidth]{media/papersWcode_tfUsage2.png}
	\end{center}
	\caption[Historical Usage of ViT in Image Tasks]{As of 2022, the usage of a Vision Transformer (ViT) in image 
	tasks matches the usage of ResNets and has outnumbered any other popular CNN architecture.
	\citep{PapersOverTime}}
	\end{figure}

Figure 1.1 was produced by PapersWithCode, a popular academic research aggregator. For the past three years, 
ResNets, the most popular architecture in image processing and computer vision, has dominated the proportion
of academic research in object detection. In 2022, Vision Transformer research popularity has reached
that of ResNets and exceeded any other major category.


In the zeitgeist of Vision Transformer research, this thesis explores a ViT-assisted metric related to 
image compression. This metric can provide additional insights to GAN output quality 
and the latent space (contextual) preservation of a variety of input images.


\NOTE{Thus, contributions from this thesis can be viewed as providing a stepping stone 
towards an end-to-end Transformer-based image compression and reconstruction framework.}


\section{Brief History}

\subsection{Attention and Language Models}

"Attention Is All You Need" is a seminal research publication
by a team of Google researchers, which kickstarted the Transformer revolution in Deep Learning in 2017.
It proposes a novel architecture, which models long-range dependencies in
sequential (text) data, by arranging a set of self-attention layers. 

A self-attention layer is what the model uses to focus on different elements of the 
input sequence simultaneously. For example, it can be used to compute the 
distance (relationship) between every word in a given sentence. \citep{Attention}


Examples of implementations of text-based Transformers are BERT by Google and GPT-3 

by OpenAI.
BERT, among many other applications, processes and autofills every single English-based 
Google user search query as of 2021 . \citep{bert}
GPT-3, on the other hand, revolutionized text generation in 2020, demonstrating the ability 
to generate extremely cohesive textual output.

Most Transformers are used for applications in language modeling and Natural Language 

Processing (NLP).
Thus, they are often benchmarked against Recurrent Neural Networks (RNNs, and specifically 
Long Short-Term Memory, LSTM architecture). LSTMs rely on hidden states to pass information 
along sequentially during the encoding and decoding process for each word token. 
However, they typically fall short learning long-range dependencies.


\subsection{Attention in Vision Tasks}

The attention mechanism is capable of focusing on objects found anywhere on an input image.
It operates within a single network layer compared to Convolutional Neural Networks (CNNs),
where the variable size convolution kernels scan across the different layers of the architecture.
\citep{dosovitskiy2020vit}

As shown in Figure 1.2 below, tokenization happens at the pixel level, i.e. each pixel attends to each other pixel in the grid. 
This becomes computationally intensive, on the order of $(n^2)^2$, where $n$ denotes width of a square image. 
To resolve this, the input image is broken down into square blocks of equal size, referred to as image patches.
Then, each image patch is unrolled into a one-dimensional sequence $(n x 1)$ and indexed with a positional 
embedding in a table for future reference and retrieval purposes. The embeddings enter the Transformer and 
finally, a feed forward classifier, in the form of a Multilayer Perceptron (MLP) 
makes the classification prediction, yielding a probability distribution. \citep{dosovitskiy2020vit}

\begin{figure}[H]
	\begin{center}
	\includegraphics[width=0.8\textwidth]{media/vit_arch.png}
	\end{center}
	\caption[ViT Architecture]{The Vision Transformer (ViT) architecture. \citep{GoogleViT}}
	\end{figure}

A Transformer, in a way, is a generalization of a feed forward network, but 
instead of fixed connections weights in an MLP, each connection weight (i.e. attention) 
is computed ad hoc. This makes the Transformer, unlike the MLP, permutation invariant. 
That is, it would not know where certain information is coming from, unless there are additional 
learnable sequential positional embeddings, i.e. index the image patches.

\section{Principles of Operation}

Continuing from the previous section, a ViT can be thought of as a 
generalization of an MLP, which itself is a generalization of a CNN. 
The ViT happens to learn very similarly to a CNN, which represents the latent space
as filters carrying principal components.

In principle, CNNs have good inductive priors and can learn any function. 
However, they promote locality, i.e. nearing pixels are probability-wise considered most important.
This may easily not be desired, especially in the key applications of object detection and, in the 
future, image compression.


The encoding process indexes embeddings. For instance, certain key words in a sentence or 
objects in image blocks are mapped in a reference lookup table. 
The Decoder outputs Keys at each step. These vectors represent hidden states,
which are being passed on into each next iteration of the Transformer. 
The last layer, expectedly, uses a Softmax architecture to normalize and map the potential output classes 
to a probability distribution. 

%\vspace{10mm}
\newpage

\textbf{Multi-Head Attention}


As shown in Figure 1.3 below, sets of parallel attention layers at each token are called multi-head attention.
This approach varies what to pay attention to, for example,at the different objects 
in an image (or in natural langue, different verbs in a sentence). 
The multi-head attention is composed of Key-Value pairs coming from 
the encoding part of the source image (i.e. the input embedding) and 
Queries from the output embedding (i.e. encoding part of target image).

\begin{figure}[H]
	\begin{center}
	\includegraphics[width=0.5\textwidth]{media/attn_diag.png}
	\end{center}
	\caption[Attention Mechanism]{The Attention mechanism architecture. \citep{GoogleViT}}
	\end{figure}


\section{Mathematical Formulation}

%\textbf{Attention}
\subsection{Attention}

In its full formulation, Attention is a function of vectors representing Queries, Keys, and Values, 
labeled as $(Q,K,V)$. 
Attention equals the dot product $(QK^T)$ of Keys and Queries respectively, softmaxed over 
the square root of dimensions and multiplied by the Values vector.

To provide further intuition:

Values are what is most interesting in the source image, i.e. 
attributes or structural features. In text, a Value could be important adjectives before each keyword, 
which provide emphasis in a given input sentence.
Keys, on the other hand, index (or address) those Values (e.g. name, type, weight). 
Each Key has an associated Value. 
Queries are built by the encoder of the target image and prompt the network to find closest 
available information (Key and its corresponding Value). 

\begin{center}
	$ Attention(Q, K, V) = softmax(\dfrac{QK^T}{\sqrt{d_k}} V ) $
	\end{center}


\NOTE{Thus, the overall dynamic is that a Query is pegged against a Key to locate a certain Value.}

%\vspace{10mm}

%\textbf{Softmax}
\subsection{Softmax}

The Softmax function is defined as a normalized exponential function. 
A sequence of variables is mapped into exponentials and divided by the sum of all exponentials. 
Thus, the large numbers become almost ones and small numbers near zeros. Softmax is similar to the maximum 
function, with a key difference that Softmax is differentiable. 

\begin{center}
	$ \displaystyle\sigma(Z_i) = \dfrac{e^{z_j}}{\sum_{j=1}^K e^{z_j}}$ for $i=1,...,K$ $and z=(z_1,...,z_K) \in \mathbb{R}^K $
	\end{center}

Thus, a Softmax of an inner product of each Key with Query vector normalizes to a 
probability distribution over all Values. Neural networks typically utilize a softmax in their last 
layer over all the classification labels. This yields the top classification by probability. 
Using the softmax function, a certain Key will stand out (in magnitude) vs the rest.

%\vspace{10mm}

%\textbf{Vector Similarity}
\subsection{Vector Similarity}

Vector proximity between embeddings represents the similarity between objects in images, or 
word connotations in sentences.


\begin{figure}[H]
	\begin{center}
	\includegraphics[width=0.5\textwidth]{media/dot_real.png}
	\end{center}
	\caption[Vector Representation]{Vector representation on a unit circle.}
	\end{figure}


The dot product of Keys and Queries yields an angle between both vectors (e.g. $u_i$ and $u_j$ in Figure 1.2 above) to measure 
how similarly aligned they are. In high dimensional spaces, most vectors would be orthonormal to each other
and $cos(90)=0$. But if Key and Query vectors are similar or align, they'd have a large dot product. 
The larger the similarity, the larger the dot product. 
The Query vector is computed with each Key in the surrounding vector space and softmaxed to 
select the one Key with the highest dot product.
The selected Key in space has an associated Value. In applications, that Value would correspond to 
a match with a labeled object in an image, or perhaps the next word in a generated sentence. 


A depiction of the vector space is shown in Figure 1.3 below.


\begin{figure}[H]
	\begin{center}
	\includegraphics[width=0.3\textwidth]{media/dot2_real.png}
	\end{center}
	\caption[Key/Query Vector Proximity]{Vector proximity shows the closest Key vector to a given Query vector in space.}
	\end{figure}


\section{Implementations}

This section reviews several notable ViT implementations of interest.
Most developments have been made open source. Howerver, some of the 
highest quality implementations are still closed source and operated 
under a license or payment wall.

\subsection{Open Source}

The academic research aggregator PapersWithCode lists 96 open source implementations of 
the original ViT from "An Image is Worth 16x16 Words" discussed in Chapter 2. 
The publication was made for ICLR 2021. \citep{Papers16}

The original model from the team of Google researchers, written using TensorFlow, has 72,567 stars in its GitHub repository.
The second highest implementation by Hugging Face, written using PyTorch, has 61,820 stars. 


\vspace{5mm}

\textbf{Pre-trained Model Used}

For the purpose of this thesis, a PyTorch implementation was used trained on ImageNet-21k
and fine tuned on ImageNet-1k. The ViT has 9,696 stars on its GitHub repository,
ranking as the 5th highest rated implementation. It was chosen due to its reliability,
project maturity, and interfaceability. \citep{Papers16}

\vspace{5mm}

A python pip package is also available on PyPI from the same "lucidrains" implementation,
which can be installed via: 
\begin{verbatim} "pip install pytorch_pretrained_vit" \end{verbatim}


Other pre-trained models are also availabe on GitHub and Google Colab. 
Some suffer from a lower GitHub star rating, are not well packaged (not object-oriented), or perhaps deviate too far
from the original implementation of "An Image is Worth 16x16 Words". 


\subsection{Closed Source}


Several modified implementations targeting specific generative applications have been made by Hugging Face and OpenAI.
OpenAI has expanded from its success from GPT-3, a text generation Transformer, to visual applications.
Transformers and Vision Transformers can be used to find context in image objects from the latent space of an input image.
CLIP, Image GPT, DALL.E, and DALL.E2 by OpenAI bridge the gap between textual information and images.
DALL.E 2, the highest quality image Transformer to date, was released in April 2022.
It is trained on 250M image-text pairs to be able to generate coherent images from textual description. \citep{Dalle2}

\NOTE{Most of OpenAI's Transformer developments (e.g. GPT-3, DALL.E 2) have not been open sourced. 
Google's ViT and BERT, however, have been.}


\section{Computational Constraints for Training}


The main reason for using a pre-trained ViT is that the Transformer architecture is so generalizable and 
overparameterized, that the computational requirements for Time, Memory, and Cost of training 
are beyond the scope of this thesis. 


For example, GPT-3 was trained on 45TB of text data (e.g. all of Wikipedia included). It has 175B parameters and 96 attention layers. \citep{GPT3}

One flavor of GPT-3 (of many) alone cost \$4.6M to train using the most advanced Nvidia datacenter GPUs in a cluster.

The original ViT implementation from "An Image is Worth 16x16 Words" took 2.5k core-days of training time. \citep{dosovitskiy2020vit} 

\vspace{5mm}

\textbf{Training the Next Generation of Image GPT}

It would require several orders of magnitude more resources to train a Vision Transformer that matches the quality of GPT-3 
for generative text applications.

The training data would need to be similar to all Google images on the internet (or at least a respectable representative subset).
It would probably cost on the order of \$50M in training alone, since image matrices are two-dimensional sequences of data.
Contextual complexity also rises exponentially for images compared to text.














%============================================================================%
%
%	CHAPTER 2
%
%============================================================================%
\cleardoublepage
\chapter{Specific Tools for a Thesis in Computer Science}

This chapter covers more specific examples, which are useful in the field of computer sciences.

\newpage
\section{"An Image is Worth 16x16 Words"}

\large{Transformers for Image Recognition at Scale}

Now that we explained relevant transformer components, we can see how it applies 
to 2D signals, i.e. image matrices for classification purposes in image recognition.
The Vision Transformer or ViT, was very recently published for ICLR 2021 by a 
google team for a ImageNet trained transformer.

The attention mechanism is capable of focusing on objects anywhere on the image 
within a single network layer versus the variable size convolution kernels 
across the different layers.

"image is worth 16x16 words" completely discards the notion of convolutions. 
Compared to a convolutional counterpart (the ResNet), the ViT cost 75\% 
less to train and beats accuracy on ImageNet by 1\%. 
ViT uses approximately 2-4x less compute to attain the same performance 
(averaged over 5 datasets). Unlike variable size convolution kernels accross 
layers, the block size in the image transformer is able to pay attention within 
a single layer to anywhere on the image. 
CNNs have good inductive priors and can learn any function. 
However, this promotes locality (i.e. nearing pixels are probably most important).
One way to think of a ViT is a generalization of an MLP which itself is a 
generalization of a CNN. The ViT also learns very similarly to a CNN 
(e.g. filters with principal components).
Transformer, in a way, is a generalization of a feed forward network, but 
instead of fixed connections weights in an MLP, each connection weight (i.e. attention) 
is computed on the fly. That makes the Transformer, unlike the MLP, permutation invariant. 
In that it wouldn't know where information is coming from unless there are additional 
learnable sequential positional embeddings, i.e. number down the image patches.

% \begin{figure}[H]
% \begin{algorithm}[H]
% \KwData{A as Array to sort,}
% \KwResult{A as sorted Array} 

% int n $\leftarrow$ A.size \tcp*[l]{cache the initial size of A}


% \Repeat{n>1}{
% 	int newn $\leftarrow$ 1
% 	\For{int i=0; i<n-1; i++}
% 	{
% 		\If{A[i] > A[i+1]}
% 		{
% 			A.swap(i, +1);
% 		}
% 	}
% 	n  $\leftarrow$ newn
% }

% return A
% \caption{bubbleSort(Array A)}
% \end{algorithm}
%\caption[TBC]{Figure caption.v.}
% \end{figure}


\newpage
\section{"Towards End-to-End Image Compression with Transformers"}


% \begin{table}[H]
% \begin{center}
% \begin{tabular}{|c|c|c|c|c|c|c|}
% \hline
% result XY & \multicolumn{6}{c|}{Rater A}\\\hline
%  \multirow{6}{*}{Rater B}& & \textbf{9} & \textbf{0} & \textbf{1} & \textbf{2} & \textbf{c(a) }\\\cline{2-7}
%  & \textbf{9}& 27	&0&	0	&0 & 0.5510204082\\\cline{2-7}
% &\textbf{0}	&0	&6	&0	&0&0.1224489796\\\cline{2-7}
% &\textbf{1}	&0	&1	&7	&0&0.1632653061\\\cline{2-7}
% &\textbf{2}	&0	&1	&7	&0&0.1632653061\\\cline{2-7}
% &\textbf{c(r)} & 0.5510204082 & 0.1632653061 & 0.2857142857 & 0& n=49\\\hline

% \end{tabular}
% \caption[Example data for mathematical equations]{A set of example data for further use in a mathematical equation.  }
% \end{center}
% \end{table}


\section{Image Generation with GANs}

Generative models.
GANs used to be the coolest thing until Transformers came along, but nonetheless 
still great. Typically, GANs pit 2 CNNs against each other, one called Generator to 
generate an image and one called Discriminator to determine how fake the image is 
compared to the input, relaying feedback to improve the Generator after each epoch. 
A GAN pits two Neural Networks (e.g. Convolutional) one against the other to generate 
new content (e.g. an image). The structure of a GAN is foundational to its design. 
As Fig.1 shows, the Generator CNN is shown on the left of the generated output. 
The Generator iteratively learns to create new synthetic data resembling real source data 
(e.g. regular or 360° images).  On the right of the generated output is the Discriminator CNN.
The Discriminator is trained to differentiate between real and synthetic data. 
In essence, the Discriminator gives the Generator feedback on its performance while 
simultaneously achieving incremental improvement in discernibility. 
Thus, each epoch the Generator produces more and more realistic data and the Discriminator iteratively improves its own ability to differentiate between real and synthetic data.


\section{First Principles of Neural Image Compression}

Deterministic, Probabilistic
GANs
JPEG/MPEG

\section{Commonly Used Metrics for Image Quality}
Metrics for the generated output image are: a BRISQUE score (reference-less) [5], 
Mean Squared Error (MSE) and Structural Similarity Index (SSIM). PSNR.
GANs use FID and Inception score. 
BRISQUE \citep{BRISQUE}
Divisive normalization, Gaussian properties from natural scene statistics
Process results in Gaussian unless it's distorted
BRISQUE too simple for user-generated content (UGC)

%============================================================================%
%
%	BIBLIOGRAPHY
%
%============================================================================%

\cleardoublepage
\phantomsection %hyperref package support
\addcontentsline{toc}{chapter}{Bibliography} % add entry to table of contents
\pagestyle{plain}


%{\textbf{\LARGE{Bibliography}}}\\	%headline
%\nocitep{*} % Show all Bib-entries (DEBUG)

\bibliographystyle{abbrvnat}
% \bibliography{bib/algorithm.bib,bib/resDes.bib} % for a better structure you can split your bib items into seperate files
\bibliography{bib/resDes.bib}


\end{document}